\section{Preliminaries}
\label{sec-pre}
Let $\Gamma$ be a countably 
infinite set of labels.

\stitle{Graphs}. 
Consider directed labeled graphs $G=(V, E, L)$,
where (1) $V$ is a finite set of vertices;
(2) $E\subseteq V\times \Gamma \times V $ is a finite set 
of edges, 
where $\langle v_1, r, v_2\rangle$ is an edge
from $v_1$ to $v_2$,  
\revise{carrying} label $r\in \Gamma$;
and (3) $L$ is a labeling function that 
assigns a label $L(v) \in \Gamma$
to each vertex $v$ in $V$.
Let $R_G$ be the set of all edge labels in $G$.
\looseness=-1

For each vertex $u {\in} V$, 
denote by $N^+_r(u)$ (resp. $N^-_r(u)$) the set of 
its  \revise{{\em outgoing} (resp. {\em incoming}) neighbors} $u_1$ of $u$
\revise{such that there exists} an  
edge labeled $r$ from $u$ to $u_1$
(resp.  
from $u_1$ to $u$).
That is, $N^+_r(u) {=} \{u_1\mid \langle u, r, u_1\rangle {\in} E\}$
and $N^-_r(u) {=} \{u_1\mid \langle u_1, r, u\rangle {\in} E\}$.
Denote by $N(u)$ the set of all neighbors of $u$, 
\ie $N(u) {=} \bigcup_{r\in R_G} (N^-_r(u) \cup N^+_r(u))$.

 

\stitle{Isomorphism and homomorphism.} 
We introduce the semantics of 
subgraph
isomorphism and homomorphism.
\revise{Consider a pattern  
$Q{=} (V_Q, E_Q, L_Q)$} and a graph $G {=} (V_G, E_G, L_G)$.
Here, a pattern is also a directed 
labeled graph, introduced
to simplify the presentation.  
\looseness=-1

 
\etitle{Subgraph isomorphism}.
An \textit{isomorphic mapping} from  
\revise{$Q$ to $G$},
denoted by 
\revise{$Q \sqsubseteq G$} is an injective mapping $\varphi$
from $V_Q$ to $V_G$  
such that
(1) for each $u{\in} V_Q$, $u$ and $\varphi(u)$ 
carry the same label, \ie $L_Q(u) = L_G(\varphi(u))$; and
(2) the adjacency relation \revise{preserves},  
\ie~\revise{$\langle u, r, u'\rangle \in E_Q$ if and only 
	if $\langle \varphi(u), r, \varphi(u')\rangle \in E_G$}. 
When such an injective $\varphi$ exists, 
we say that $Q$ is subgraph isomorphic to $G$. 
 
 